\section{Segre coordinates in the same entire original function field}
\label{sec:projective-segre-native-coordinates}
\begin{lemma}[The actual generic point retains every original chart coefficient]
\label{lem:segre-native-generic-chart-coefficients}
\lean{EnvelopeProperty.ProjectiveSegreNativeCoordinates.nativeChart,EnvelopeProperty.ProjectiveSegreNativeCoordinates.nativeChart_fac,EnvelopeProperty.ProjectiveSegreNativeCoordinates.nativeChart_coordinate}
\leanok
\uses{thm:projective-chart-coefficient-algebra,lem:fc-original-projective-chart-section-ratios}
Let $X$ be integral over a field $k$, and let $f:X\to\mathbf P^m_k$ be any
morphism. If the inverse image of the $i$th standard chart is nonempty,
the actual generic-stalk map constructs a morphism
$\operatorname{Spec}K(X)\to\operatorname{Spec}k[X_a/X_i]$ whose composite
with the chart inclusion is the original generic point followed by $f$.
Every polynomial coordinate pulls back to the same original rational
function $f^*(X_a/X_i)\in K(X)$.
\end{lemma}
\begin{proof}
The original generic point lies in every nonempty open. Apply the original
stalk homomorphism of $f$, the canonical morphism from the spectrum of the
image stalk into that open, and its actual polynomial-chart isomorphism.
The canonical stalk square and the open-stalk factorization give the
entire asserted morphism triangle. For a coordinate coefficient, the
canonical open-to-spectrum map composed with the stalk morphism is the
spectrum of the original germ homomorphism. Consequently taking global
sections and the canonical global-sections identification on a spectrum
gives exactly the original germ. The original chart-coordinate formula
identifies that coefficient with $X_a/X_i$.
\end{proof}
\begin{theorem}[Every Segre coordinate is the product of the original factor ratios]
\label{thm:segre-native-coordinate-product}
\lean{EnvelopeProperty.ProjectiveSegreNativeCoordinates.nativeProductChart,EnvelopeProperty.ProjectiveSegreNativeCoordinates.nativeProductChart_fac,EnvelopeProperty.ProjectiveSegreNativeCoordinates.segre_chart_nonempty,EnvelopeProperty.ProjectiveSegreNativeCoordinates.coordinate_product}
\leanok
\uses{lem:segre-native-generic-chart-coefficients,thm:original-segre-pulled-coordinate-section}
Let $f:X\to\mathbf P^m_k$ and $g:X\to\mathbf P^n_k$ have the same original
base triangle. Suppose their $i$th and $j$th inverse-image charts are
nonempty. For the actual original pair followed by the Segre immersion,
the $(i,j)$th inverse-image chart is nonempty and every coordinate ratio
in the same original $K(X)$ satisfies
\[
 (f,g)^*\frac{Z_{a,b}}{Z_{i,j}}
 = f^*\frac{X_a}{X_i}\,g^*\frac{Y_b}{Y_j}.
\]
No closed-immersion, characteristic, perfectness or coordinate-comparison
hypothesis is imposed.
\end{theorem}
\begin{proof}
The two actual generic chart maps commute over the original coefficient
field because their projective composites are the original generic point
followed by $f$ and $g$. The genuine tensor-chart pullback therefore
constructs their unique pair map. Its composite with the actual product
chart inclusion agrees with the original generic point followed by the
original pair: both projective projections agree. Compose with the Segre
immersion and use its actual chart formula. Evaluation at the closed point
of $\operatorname{Spec}K(X)$ proves that the original generic point lies
in the required Segre chart, so this chart is nonempty. Cancellation of
its open immersion identifies the constructed affine composite with the
canonical generic Segre chart map. On the affine tensor chart every
coordinate pulls back to the product of its two factor coordinates.
Taking the canonical global-sections identification of $\operatorname{Spec}K(X)$
and applying the preceding lemma gives the displayed equality in the
entire original function field.
\end{proof}
\begin{theorem}[One original closed factor retains the entire original function field]
\label{thm:segre-native-entire-function-field}
\lean{EnvelopeProperty.ProjectiveSegreNativeCoordinates.adjoin_coordinates_eq_top}
\leanok
\uses{thm:segre-native-coordinate-product,thm:projective-chart-function-field-generation}
If the original $f$ is a closed immersion, the actual Segre coordinate
ratios generate the entire same original field $K(X)$ over $k$.
\end{theorem}
\begin{proof}
For each $a$, the Segre ratio indexed by $(a,j)$ is the original
$f^*(X_a/X_i)$ times the original pivot ratio $g^*(Y_j/Y_j)=1$.
Thus the field generated by the actual Segre ratios contains every
original ratio of $f$. The genuine closed-immersion function-field
generation theorem says that these original ratios generate all of
$K(X)$. The field generated by the Segre ratios is therefore the whole
field as well.
\end{proof}
\begin{theorem}[The original Segre composite is a genuine closed embedding]
\label{thm:original-segre-composite-closed-embedding}
\lean{EnvelopeProperty.ProjectiveSegreNativeCoordinates.morphism_isClosedImmersion}
\leanok
\uses{def:original-pair-segre-morphism}
If $f$ is a closed immersion, the actual original Segre composite is a
closed immersion as well. This assertion does not require integrality.
\end{theorem}
\begin{proof}
The first projection from the original projective fiber product is
separated because it is a base change of the separated original second
projective space. The pair morphism followed by this projection is
literally $f$, a closed immersion. Cancellation of a separated target
morphism therefore makes the pair morphism a closed immersion. The
actual Segre map is a closed immersion, and composing these two closed
immersions proves the assertion.
\end{proof}
